\chapter{Wiring Diagrams}

This chapter shows basic wiring examples for common projects using the ESP32 RC Controller firmware.

These diagrams are simplified so beginners can quickly understand how to connect the hardware.

\section{Basic Robot Wiring}

The most common project is a simple differential drive robot using two motors.

\begin{center}
	\begin{tikzpicture}[
		block/.style={draw, rectangle, minimum width=3cm, minimum height=1cm},
		line/.style={->, thick}
		]

		\node[block] (phone) {Android App};
		\node[block, right=3cm of phone] (esp32) {ESP32};
		\node[block, right=3cm of esp32] (driver) {Motor Driver};
		\node[block, above right=1cm and 1.5cm of driver] (motor1) {Left Motor};
		\node[block, below right=1cm and 1.5cm of driver] (motor2) {Right Motor};

		\draw[line] (phone) -- node[above]{Bluetooth} (esp32);
		\draw[line] (esp32) -- node[above]{Control Signals} (driver);
		\draw[line] (driver) -- (motor1);
		\draw[line] (driver) -- (motor2);

	\end{tikzpicture}
\end{center}

\section{Servo Control Example}

Servos can be connected directly to PWM pins on the ESP32.

\begin{itemize}
	\item Servo signal → ESP32 GPIO pin
	\item Servo power → 5V supply
	\item Servo ground → common ground
\end{itemize}

Example uses:

\begin{itemize}
	\item Steering control
	\item Camera pan control
	\item Robotic arm joints
\end{itemize}

\section{LED and Buzzer Outputs}

Digital outputs can control LEDs or buzzers.

Example connections:

\begin{itemize}
	\item GPIO pin → LED (with resistor)
	\item GPIO pin → buzzer module
\end{itemize}

These outputs can be controlled using buttons or switches in the Android application.

\section{Sensor Connections}

Sensors can be connected using:

\begin{itemize}
	\item Analog inputs
	\item Digital inputs
	\item I2C interface
\end{itemize}

The firmware allows sensor data to be transmitted to the Android application using telemetry.